% ===== FILE: sources/2503.01840-eagle-3-scaling-up/arxiv/paper.tex =====
% This must be in the first 5 lines to tell arXiv to use pdfLaTeX, which is strongly recommended.
\pdfoutput=1
% In particular, the hyperref package requires pdfLaTeX in order to break URLs across lines.

\documentclass[11pt]{article}

% Change "review" to "final" to generate the final (sometimes called camera-ready) version.
% Change to "preprint" to generate a non-anonymous version with page numbers.
\usepackage[final]{acl}

% Standard package includes
\usepackage{times}
\usepackage{latexsym}

% For proper rendering and hyphenation of words containing Latin characters (including in bib files)
\usepackage[T1]{fontenc}
% For Vietnamese characters
% \usepackage[T5]{fontenc}
% See https://www.latex-project.org/help/documentation/encguide.pdf for other character sets

% This assumes your files are encoded as UTF8
\usepackage[utf8]{inputenc}

% This is not strictly necessary, and may be commented out,
% but it will improve the layout of the manuscript,
% and will typically save some space.
\usepackage{microtype}

% This is also not strictly necessary, and may be commented out.
% However, it will improve the aesthetics of text in
% the typewriter font.
\usepackage{inconsolata}

%Including images in your LaTeX document requires adding
%additional package(s)
\usepackage{graphicx}


\usepackage{subcaption}
\usepackage{amsmath}
\usepackage{float}
\usepackage{multirow}
\usepackage{booktabs}
\usepackage{url}
\usepackage{xcolor}
\definecolor{skyblue}{RGB}{135,206,235}


\title{\includegraphics[width=0.06\textwidth]{./figs/logo.png} EAGLE-3: Scaling up Inference Acceleration of Large Language Models via Training-Time Test}



\author{Yuhui Li\textsuperscript{1,3}, Fangyun 
Wei\textsuperscript{2}, Chao Zhang\textsuperscript{1}, Hongyang 
Zhang\textsuperscript{3,4}\\
 \textsuperscript{1}Peking University 
 \textsuperscript{2}Microsoft Research \\
 \textsuperscript{3}University of Waterloo
 \textsuperscript{4}Vector Institute\\
 \texttt{yuhui.li@stu.pku.edu.cn, fawe@microsoft.com}\\
 \texttt{c.zhang@pku.edu.cn, hongyang.zhang@uwaterloo.ca}
 }


\begin{document}
\maketitle






\begin{abstract}

The sequential nature of modern LLMs makes them expensive and slow, and speculative sampling has proven to be an effective solution to this problem. Methods like EAGLE perform autoregression at the feature level, reusing top-layer features from the target model to achieve better results than vanilla speculative sampling. A growing trend in the LLM community is scaling up training data to improve model intelligence without increasing inference costs. However, we observe that scaling up data provides limited improvements for EAGLE. We identify that this limitation arises from EAGLE's feature prediction constraints. In this paper, we introduce EAGLE-3, which abandons feature prediction in favor of direct token prediction and replaces reliance on top-layer features with multi-layer feature fusion via a technique named training-time test. These improvements significantly enhance performance and enable the draft model to fully benefit from scaling up training data. Our experiments include both chat models and reasoning models, evaluated on five tasks. The results show that EAGLE-3 achieves a speedup ratio up to 6.5x, with about 1.4x improvement over EAGLE-2. In the SGLang framework, EAGLE-3 achieves a 1.38x throughput improvement at a batch size of 64. The code is available at \url{https://github.com/SafeAILab/EAGLE}.
\end{abstract}

\begin{figure}

  \begin{subfigure}{\linewidth}
    \centering
    \includegraphics[width=\linewidth]{figs/scaale.pdf}
    %\caption{Speedup ratio.}
    \label{fig:sc02}
  \end{subfigure} \hfill
  \begin{subfigure}{\linewidth}
    \centering
    \includegraphics[width=\linewidth]{figs/scaletau.pdf}
    %\caption{Average acceptance length.}
    \label{fig:sc03}
  \end{subfigure}
  \vspace{-0.8cm}
  \caption{Scaling law evaluated on the MT-bench using LLaMA-Instruct 3.1 8B as the target model, with the x-axis representing the data scale relative to ShareGPT. The new architectural designs in EAGLE-3 enable an increasing scaling curve, which was never observed in the previous works.}
  \label{fig:sc}
  \vspace{-0.5cm}
\end{figure}

\section{Introduction}







\begin{figure*}[t]
\centering
  \includegraphics[width=\textwidth]{figs/r1.pdf}
  \vspace{-0.3cm}
  \caption{Speedup ratios of different methods at temperature=0. For the standard speculative sampling, Vicuna-13B uses Vicuna-68M as the draft model. In Table \ref{tab:experiments}, we present comparisons with additional methods, but this figure only showcases a subset. Chat model's evaluation dataset is MT-bench, and the reasoning model's evaluation dataset is GSM8K. DeepSeek R1 LLaMA 8B refers to DeepSeek-R1-Distill-LLaMA 8B.}
  \label{fig:speedup}
  \vspace{-0.3cm}
\end{figure*}

Modern Large Language Models (LLMs) are being applied to more domains, with their improved capabilities driven by scaling model parameters—some LLMs now exceed hundreds of billions of parameters. In autoregressive generation, each token requires accessing all model parameters, making LLM inference slow and costly.

Recently, test-time scaling up has gained significant attention. Models like ChatGPT o1 and DeepSeek-R1~\cite{guo2025deepseek} engage in deliberate reasoning before responding, pushing the boundaries of LLM capabilities at the cost of longer inference time. However, these models often require lengthy reasoning processes, making them extremely costly, while the increased response time severely impacts user satisfaction. These reasoning models significantly increase the proportion of inference costs in the overall LLM pipeline, driving researchers to explore cheaper and faster inference optimization methods.



Speculative sampling methods can reduce LLM latency by partially parallelizing the generation process. These methods rapidly generate draft tokens and then verify them in parallel. This allows multiple tokens to be produced in a single forward pass, significantly reducing inference latency. In vanilla speculative sampling, the draft model is a separate, smaller LLM, typically a lower-parameter version from the same series as the target model. This draft model operates independently of the target model. Unlike the vanilla speculative sampling, EAGLE~\cite{li2024eagle} reuses the top-layer features of the target model (the features before the LM head). It trains the draft model to autoregressively predict the next feature and then uses the target model's LM head to obtain the draft token. By leveraging the rich information from the target model, EAGLE achieves significantly better acceleration compared to vanilla speculative sampling. Subsequent methods such as HASS~\cite{zhang2024learning} and Falcon~\cite{gao2024falcon} also adopt the approach of predicting the next feature using the current feature sequence.

Recent LLMs have increasingly relied on larger training datasets to achieve better performance. For example, LLaMA series models with sizes of 7B (8B) have used 1T, 2T, and 15T tokens of training data for LLaMA 1~\cite{touvron2023llama}, LLaMA 2~\cite{touvron2023llama2}, and LLaMA 3~\cite{dubey2024llama}, respectively, resulting in significant improvements across various metrics while keeping the model architecture and inference cost largely unchanged. Similarly, we aim to improve the acceptance rate and acceleration ratio of EAGLE by increasing its training data. Unfortunately, we observe that the gains from additional training data for EAGLE are limited. 
We analyze the reasons behind this phenomenon. As shown in the upper part of Figure \ref{fig:nofe}, EAGLE performs autoregressive prediction at the feature level, predicting the next feature and then feeding the feature into the LM head of the target model to obtain the token distribution. EAGLE's loss function consists of two components: the feature prediction loss $l_{\text{fea}}$ and the token prediction loss $l_{\text{token}}$. Thanks to the feature prediction loss, the draft model trained only at Step 1 can adapt to Step 2 and acquire multi-step prediction capabilities. However, with token prediction as the ultimate goal, feature prediction can be seen as an additional constraint, which limits the expressiveness of the draft model and makes it difficult to benefit from increased data. After removing the feature constraint and expanding the training data (the middle part of Figure \ref{fig:nofe}), as shown in Figure \ref{fig:sc1}, the acceptance rate $0$-$\alpha$ of the first draft token improves significantly.
However, the output of the draft model in Step 1, denoted as $\hat{a}_{t+1}$, is far away from the ground-truth $f_{t+1}$, causing the input sequence $f_1,f_2,\cdots,f_t,\hat{a}_{t+1}$ in Step 2 to deviate significantly from the training distribution, resulting in a very low acceptance rate $1$-$\alpha$ for the second draft token, as shown in Figure \ref{fig:sc1}. We can address this issue by incorporating Step 1 into the training process (the bottom of Figure \ref{fig:nofe}). Using this method, the benefits of increasing training data become more pronounced. We name this technique as training-time test.

\begin{figure}[t]
  \centering
  \includegraphics[width=1.07\columnwidth]{figs/nofe.pdf}
  \vspace{-0.5cm}
  \caption{Illustration of \textbf{training-time test} (the bottom part) and its comparison with other draft methods (the upper and middle parts). $f$ denotes the feature, $t$ denotes the token, and $a$ represents the unconstrained vectors. We use the hat to denote the predictions from models. All the methods shown in the figure use the token sequence from the previous time step, but for simplicity, this is not depicted in the figure. The input to EAGLE-3 is not actually $f$, but it is not shown in this figure. We will provide a detailed explanation in the following section.}
  \label{fig:nofe}
  \vspace{-0.5cm}
\end{figure}

\begin{figure}

  \begin{subfigure}{0.49\linewidth}
    \centering
    \includegraphics[width=\linewidth]{figs/alpha0.pdf}
    %\caption{Speedup ratio.}
    \label{fig:sc00}
  \end{subfigure} \hfill
  \begin{subfigure}{0.49\linewidth}
    \centering
    \includegraphics[width=\linewidth]{figs/alpha1.pdf}
    %\caption{Average acceptance length.}
    \label{fig:sc01}
  \end{subfigure}
  
  \caption{Comparison of acceptance rates across different methods, with the x-axis representing the data scale relative to ShareGPT.}
  \label{fig:sc1}
  \vspace{-0.5cm}
\end{figure}


EAGLE and speculative sampling methods such as Medusa~\cite{cai2024medusa} reuse the top-layer features of the target model, specifically the features immediately before the LM head. For an LM head with a full-rank weight matrix, the top-layer features corresponding to the logits of the next token are unique, ensuring that the information contained in these features aligns directly with the logits of the next token. However, predicting the next-next token based solely on top-layer features—which are inherently limited to the next token—poses a significant challenge. Fortunately, the training-time test technique described above enables the use of features from intermediate layers instead of relying solely on the top layer, as the feature prediction loss $l_{\text{fea}}$ has been removed during training.

To summarize, this paper introduces EAGLE-3, an enhanced version of EAGLE that achieves a significant speedup. EAGLE-3 is parallelized and fully compatible with the drafting tree technique from EAGLE-2~\cite{li2024eagle2}. Our key contributions include:
\begin{itemize}
    \item 
    \textbf{A novel training-time test architecture for the draft model:} We remove the feature prediction constraint and directly predict tokens while simulating multi-step generation during training. This direct token prediction provides complete flexibility in the draft model’s input. Instead of reusing only the top-layer features, we integrate and leverage low-, mid-, and high-level features from the target model, capturing rich semantic information from different layers.
    \item
    \textbf{Discovery of a scaling law for inference acceleration in large language models:} With the new architecture, we observe that increasing the amount of training data for the draft model leads to a proportional increase in the speedup ratio of EAGLE-3. This scaling behavior was not observed in the original EAGLE architecture, as shown in Figure \ref{fig:sc}
    \item
    \textbf{Improved inference acceleration:} EAGLE-3, trained with approximately 8x more data than EAGLE, achieves a 1.4x latency speedup over EAGLE-2 at batch size 1. Speculative sampling is often thought to reduce throughput at large batch sizes. However, in SGLang~\cite{zheng2024sglang}, a production-grade framework, EAGLE-3 improves throughput by 40\% at a batch size of 64. We expect larger data size would lead to further improved speedup ratio.
\end{itemize}



\section{Preliminaries}
\subsection{Speculative Sampling}
Speculative sampling \cite{leviathan2023fast,chen2023accelerating,sun2024spectr,sun2024optimal} is a lossless LLM acceleration technique that alternates between drafting and verification, where drafting is performed at low cost and verification is parallelized, corresponding to the generation of drafts and the verification process, respectively. We use $t_i$ to denote the $i$-th token and $T_{a:b}$ to represent the token sequence $t_a, t_{a+1}, \cdots, t_b$. When $T_{1:j}$ is used as the prefix, the two stages of speculative sampling are as follows.

In the drafting stage, speculative sampling utilizes a draft model (a smaller version from the same series as the target model) to autoregressively generate $k$ tokens to form the draft. $\hat{T}_{j+1:j+k}$, while also recording the probability $\hat{p}$ for each token.

In the verification stage, speculative sampling invokes the target model to evaluate the draft $\hat{T}_{j+1:j+k}$ and records its probability $p$. Speculative sampling then determines the acceptance of draft tokens sequentially, from front to back. For token $\hat{t}_{j+i}$, the probability of acceptance is given by $\min(1, p_{j+i}(\hat{t}_{j+i})/\hat{p}_{j+i}(\hat{t}_{j+i}))$. If the token is accepted, the process moves to the next token. Otherwise, a token is sampled from the distribution $\text{norm}(\max(0, p_{j+i} - \hat{p}_{j+i}))$ to replace $\hat{t}_{j+i}$, and the remaining tokens in the draft are discarded. 
Appendix A.1 of \cite{leviathan2023fast} proves that speculative sampling is consistent with the distribution of vanilla autoregressive decoding. 
% EAGLE, EAGLE-2, and EAGLE-3 adhere to this framework, ensuring that the output distribution of the inference pipeline provably matches that of the target model. 
% Throughout this paper, we neither change the target model nor relax the acceptance conditions.
\subsection{EAGLE and EAGLE-2}

The draft model with limited capacity struggles to precisely approximate the large-scale target model. EAGLE leverages the top-layer features of the target model as additional information and performs autoregression at the feature level, simplifying the drafting process. EAGLE performs autoregression at the feature level and then uses the LM head of the target model to obtain the draft token. Due to the sampling results at the token layer being hidden, feature-level autoregression introduces uncertainty. EAGLE addresses this issue by feeding the token sequence from the previous time step, i.e., the sampling results, into the draft model.
Unlike the chain-like drafts of Vanilla speculative sampling, EAGLE generates multiple draft tokens at the same position, resulting in a tree-like draft. In the verification stage, EAGLE uses tree attention to parallelize the verification of the draft tree. Interestingly, EAGLE inspired the \emph{multi-token prediction} technique used in the pre-training of DeepSeek-v3~\cite{liu2024deepseek}, which in turn inspired new architectural designs in EAGLE-3.

EAGLE~\cite{li2024eagle} and Medusa~\cite{cai2024medusa}, among others, use tree-shaped drafts, where the structure of the draft tree is predefined, static, and context-independent. The difficulty of drafting is closely related to the context, and a static draft tree can lead to resource wastage. EAGLE-2~\cite{li2024eagle2} approximates the acceptance rate using the confidence of the draft model and dynamically generates the draft tree based on this, performing pruning of the draft tree at the end of the drafting stage. EAGLE-3 also adopts the context-aware dynamic draft tree proposed in EAGLE-2.


\section{EAGLE-3}
In this section, we provide a detailed description of the implementation of EAGLE-3.

\subsection{Inference Pipeline}
\label{sec:inf pip}

Consistent with other speculative sampling methods, EAGLE-3 alternates between the drafting and verification stages. The difference between EAGLE-3 and EAGLE lies in the drafting stage, which we introduce with an example, as shown in Figure \ref{fig:pipline}. Consider the prefix ``How can''.
During the prefill phase or the previous verification stage, the target model performs a forward pass to generate the next token, ``I''.
We record the low, middle, and high-level feature sequences from the target model’s forward pass, denoted as $l$, $m$, and $h$, respectively. We concatenate the $k$-dimensional vectors $l$, $m$, and $h$ to form a $3k$-dimensional vector, then pass it through a fully connected (FC) layer to reduce it to $k$-dimensions, obtaining a feature $g$ that integrates information from different layers.
Here, $k$ refers to the hidden size of the target model.

\begin{figure}[t]
  \includegraphics[width=\columnwidth]{figs/method.pdf}
  \caption{Diagram of the EAGLE-3 inference pipeline, illustrating the three steps of the draft model. $l$, $m$, and $h$ represent the low, middle, and high-level features of the target model, respectively. $e$ denotes the embedding.}
  \label{fig:pipline}
  \vspace{-0.5cm}
\end{figure}

Our goal is to generate a draft token sequence with the prefix ``How can I''. By inputting only $g_\text{how}$ and $g_\text{can}$, the draft model cannot access the random sampling process. Therefore, similar to EAGLE~\cite{li2024eagle}, we introduce the embedding $e_\text{I}$ of the sampled token ``I''. 
The concatenated vector is then passed through an FC layer to reduce its dimensionality to $k$, and subsequently inputted into a single layer decoder, producing the output $a$.
Finally, we input $a_\text{I}$ into the LM head and sample to obtain the draft token ``do''.

In Step 1, with the prefix ``How can'', we reuse $g_\text{how}$ and $g_\text{can}$ from the target model. 
In Step 2, the prefix becomes ``How can I''. Ideally, we would reuse $g_\text{how}$, $g_\text{can}$, and $g_\text{I}$ from the target model. 
However, this is not possible because the token ``I'' has not yet been checked by the target model, and we cannot obtain $g_\text{I}$.
Instead, we use the output $a_\text{I}$ from the draft model in the previous step to replace $g_\text{I}$, and concatenate $a_\text{I}$ with the embedding $e_\text{do}$ of the sampled result ``do'' as the input to the draft model in Step 1.
In Step 3, we similarly cannot obtain $g_\text{do}$, so we use $a_\text{do}$ as a replacement, concatenating $a_\text{do}$ with  $e_\text{it}$ as the input to the draft model. The same approach is followed for subsequent steps.



\subsection{Draft Model Training}

\begin{figure*}[t]
  \includegraphics[width=\textwidth]{figs/train.pdf}
  \vspace{-0.5cm}
  \caption{Diagram of the attention causal masks during training-time test. It sequentially shows a native training step (the first step) and two simulated training steps (the second and third steps). The arrows between tokens represent contextual relationships. The \textcolor{gray}{gray} tokens represent the training data while the $\color{skyblue}{\text{blue}}$ and \textcolor{yellow}{yellow} tokens represent the first- and second-round predictions by the draft model, respectively.}
  \label{fig:mask}
  % \vspace{-0.2cm}
\end{figure*}

The input to the draft model in EAGLE is either, or at least approximately, the top-layer features $f_1,f_2,\cdots,f_t$ of the target model.
In contrast, the input to the draft model in EAGLE-3 may include the features $g_1,g_2,\cdots,g_t$ from the target model, or it may include the output $a_{t+1},a_{t+2}\cdots,a_{t+j}$ from the draft model. Therefore, we need to train the draft model to adapt to different inputs. During training, we perform test steps, where we generate $a$ and feed it back into the draft model for further training.

The core of the draft model in EAGLE-3 is a Transformer decoder layer. Aside from the self-attention operation, no other components interact with the context, so no further modifications are required during training or testing. The only component that requires slight modification is the self-attention, which we will describe in detail below.



Although the actual input consists of features, for clarity, we describe the process using tokens as input. As shown in Figure \ref{fig:mask}, the original training data is a sequence of length 3, ``How can I'', with a normal sequential dependency in the context. Therefore, the attention mask is a standard lower triangular matrix. The outputs at the three positions are ``are'', ``we'', and ``do'', which have a tree-like contextual relationship with ``how'', ``can'', and ``I''. As a result, when the input ``are'', ``we'', and ``do'' is fed into Step 2, the attention mask needs to be adjusted accordingly, as shown in the top-right corner of Figure \ref{fig:mask}. All attention masks are diagonal, except when the original training data is used as the key. Using matrix multiplication in this case would result in significant computational waste, so we can use vector dot products to calculate the attention score only for the corresponding positions.

HASS~\cite{zhang2024learning} and EAGLE-3 both make similar modifications to the attention mechanism to simulate the testing process during training, but this is not the main focus of EAGLE-3. The motivations, methods, and outcomes of the two approaches are distinctly different. The motivation behind HASS is to mitigate the error accumulation caused by inaccurate feature predictions in EAGLE. HASS still performs feature prediction, includes a feature prediction loss $l_\text{fea}$, and the input to the draft model must be the top-layer features. In contrast, the motivation behind EAGLE-3 is to remove unnecessary constraints to enhance the model's expressive power. EAGLE-3 no longer requires the draft model's output to fit the top-layer features of the target model, thus avoiding error accumulation. After removing feature prediction, the input to EAGLE-3 is completely free, and it is replaced by a fusion of features from different layers of semantic information. The removal of the feature prediction loss also enables us to discover a new scaling law for inference acceleration which was never found before. Figure \ref{fig:speedup} also shows the speedup of EAGLE-3 and HASS, with EAGLE-3 demonstrating significantly better performance.

% \subsection{Other Adjustments}


\section{Experiments}

\textbf{Models.} We conduct experiments with state-of-the-art open-source chat and reasoning models, including Vicuna 13B  \cite{vicuna2023}, LLaMA-Instruct 3.1 8B, LLaMA-Instruct 3.3 70B~\cite{dubey2024llama}, and DeepSeek-R1-Distill-LLaMA 8B~\cite{deepseek2025deepseek}. Due to the GPU constraint, we are unable to test EAGLE-3 on the 405B and 671B models.

\textbf{Tasks.} Following EAGLE \cite{li2024eagle} and Spec-Bench \cite{xia2024unlocking}, we evaluate on five common tasks, using the same weights for all tasks without fine-tuning on the respective tasks. For multi-turn conversation, code generation, mathematical reasoning, instruction following, and summarization,, we chose the MT-bench \cite{zheng2023judging}, HumanEval \cite{chen2021evaluating}, GSM8K \cite{cobbe2021training}, Alpaca \cite{alpaca}, and CNN/Daily Mail \cite{nallapati2016abstractive} datasets, respectively.


\begin{table*}[h]
  \centering
  \caption{Speedup ratios and average acceptance lengths $\tau$ of different methods. V represents Vicuna, L31 represents LLaMA-Instruct 3.1, L33 represents LLaMA-Instruct 3.3, and DSL represents DeepSeek-R1-Distill-LLaMA. SpS denotes standard speculative sampling, with its draft model being Vicuna-68M. Methods like Medusa relax acceptance conditions under non-greedy settings, which do not guarantee lossless acceleration. Therefore, we do not compare EAGLE-3 with these methods when temperature=1.}
  \resizebox{\linewidth}{!}{
    \begin{tabular}{cccccccccccccc}
    \toprule
          &       & \multicolumn{2}{c}{MT-bench} & \multicolumn{2}{c}{HumanEval} & \multicolumn{2}{c}{GSM8K} & \multicolumn{2}{c}{Alpaca} & \multicolumn{2}{c}{CNN/DM} & \multicolumn{2}{c}{Mean} \\
    \midrule
    Model & Method & Speedup & $\tau$     & Speedup & $\tau$     & Speedup & $\tau$     & Speedup & $\tau$     & Speedup & $\tau$     & Speedup & $\tau$ \\
    \midrule
    \multicolumn{14}{c}{Temperature=0} \\
    \midrule
    \multirow{8}[2]{*}{V 13B} & SpS   & 1.93x & 2.27  & 2.23x & 2.57  & 1.77x & 2.01  & 1.76x & 2.03  & 1.93x & 2.33  & 1.92x & 2.24 \\
          & PLD   & 1.58x & 1.63  & 1.85x & 1.93  & 1.68x & 1.73  & 1.16x & 1.19  & 2.42x & 2.50  & 1.74x & 1.80 \\
          & Medusa & 2.07x & 2.59  & 2.50x & 2.78  & 2.23x & 2.64  & 2.08x & 2.45  & 1.71x & 2.09  & 2.12x & 2.51 \\
          & Lookahead & 1.65x & 1.69  & 1.71x & 1.75  & 1.81x & 1.90  & 1.46x & 1.51  & 1.46x & 1.50  & 1.62x & 1.67 \\
          & Hydra & 2.88x & 3.65  & 3.28x & 3.87  & 2.93x & 3.66  & 2.86x & 3.53  & 2.05x & 2.81  & 2.80x & 3.50 \\
          & EAGLE & 3.07x & 3.98  & 3.58x & 4.39  & 3.08x & 3.97  & 3.03x & 3.95  & 2.49x & 3.52  & 3.05x & 3.96 \\
          & EAGLE-2 & 4.26x & 4.83  & 4.96x & 5.41  & 4.22x
